%%
%% spring-lecture12.tex
%% 
%% Made by Alex Nelson <pqnelson@gmail.com>
%% Login   <alex@lisp>
%% 
%% Started on  2026-04-30T12:14:05-0700
%% Last update 2026-04-30T12:14:05-0700
%% 

\lecture{}

\begin{definition}
Given a variation field $V\in\Vect{\gamma}$, the \define{First Variation}
of the energy $E_{g}$ at $\gamma$ along $V$ is defined as
\begin{equation}
\D E_{g}(\gamma)[V]=\left.\frac{\D}{\D s}E_{g}(\gamma_{s})\right|_{s=0},
\end{equation}
where $\gamma_{s}:=F(-,s)$.
\end{definition}

\begin{remark}
If $\gamma\colon[a,b]\to M$ is a smooth curve, and if
$V\in\Vect{\gamma}$ is a variation field, then the endpoints are fixed $V(a)=V(b)=0$.
Indeed,
\begin{equation}
V(t)=\frac{\partial F}{\partial s}(t,0),
\end{equation}
so
\begin{equation}
V(a)=\frac{\partial F}{\partial s}(a,0)=\frac{\partial}{\partial s}(\gamma_{s}(a))=0,
\end{equation}
since $\gamma_{s}(a)=a$ for all $s$. Similar reasoning shows $V(b)=0$.
\end{remark}

\begin{theorem}
Let $(M,g)$ be a Riemannian manifold, let $\gamma\colon[a,b]\to M$ be
a smooth curve, let $V\in\Vect{\gamma}$ be a vector field along the
curve. Then
\begin{equation}
\D E_{g}(\gamma)[V] =
-\int^{b}_{a}g_{\gamma(t)}\left(V(t),\frac{D\gamma_{s}'}{\D t}(t)\right)\,\D t
\end{equation}
for every variation field $V$ along $\gamma$.
\end{theorem}

\begin{proof}
Note that
\begin{equation}
\gamma'_{s}(t)=\frac{\partial F}{\partial t}(t,s).
\end{equation}
As the Levi-Civita connection is metric-compatible, we have the
following facts:
\begin{equation}\label{eq:math151c:lecture12:fact1}
\frac{\partial}{\partial t}(g_{\gamma_{s}(t)}(\partial_{s}F(s,t),\gamma'_{s}(t)))
= g_{\gamma_{s}(t)}\left(\frac{D}{\D t}\left(\frac{\partial F}{\partial s}\right)(s,t),\gamma'_{s}(t)\right)+g_{\gamma_{s}(t)}\left(\frac{\partial F}{\partial s}(s,t),\frac{D\gamma'_{s}}{\D t}(t)\right)
\end{equation}
(exercise: perform this  calculation!), and
\begin{equation}\label{eq:math151c:lecture12:fact2}
\frac{\partial}{\partial s}(g_{\gamma_{s}(t)}(\gamma'_{s}(t),\gamma'_{s}(t)))
=2g_{\gamma_{s}(t)}\left(\frac{D}{\D t}\left(\frac{\partial F}{\partial s}\right)(s,t)  ,\gamma'_{s}(t)\right)
\end{equation}
These are by metric-compatibility of the connection. (We have not yet
used the torsion-free property.) Moreover, as the Levi-Civita
connection is torsion-free, we have
\begin{equation}\label{eq:math151c:lecture12:fact3}
\left(\frac{D}{\D t}\frac{\partial F}{\partial s}\right)(t,s)-
\left(\frac{D}{\D s}\frac{\partial F}{\partial t}\right)(t,s)
=\D F_{(t,s)}\left(\left[\left.\frac{\partial}{\partial t}\right|_{(t,s)},\left.\frac{\partial}{\partial s}\right|_{(t,s)}\right]\right)=0.
\end{equation}
Then
\begin{subequations}
  \begin{align}
\frac{\D}{\D s}E_{g}(\gamma_{s})
&=\frac{1}{2}\frac{\D}{\D s}\int^{b}_{a}g_{\gamma_{s}(t)}(\gamma'_{s}(t),\gamma'_{s}(t))\,\D t\quad\mbox{by defn}\\
&=\frac{1}{2}\int^{b}_{a}\frac{\partial}{\partial s}g_{\gamma_{s}(t)}(\gamma'_{s}(t),\gamma'_{s}(t))\,\D t\\
&=\frac{1}{2}\int^{b}_{a}2g_{\gamma_{s}(t)}\left(\left(\frac{D}{\D s}\frac{\partial F}{\partial t}\right)(t,s),\gamma'_{s}(t)\right)\,\D t\quad\mbox{by Fact~\eqref{eq:math151c:lecture12:fact2}}\\
&=\int^{b}_{a}g_{\gamma_{s}(t)}\left(\left(\frac{D}{\D t}\frac{\partial F}{\partial s}\right)(t,s),\gamma'_{s}(t)\right)\,\D t\quad\mbox{by Fact~\eqref{eq:math151c:lecture12:fact3}}\\
&=\int^{b}_{a}\frac{\partial}{\partial t}\left(g_{\gamma_{s}(t)}\left(\frac{\partial F}{\partial s}(t,s),\gamma'_{s}(t)\right)\right)\D t\nonumber\\
&\phantom{=\quad}-\int^{b}_{a}g_{\gamma_{s}(t)}\left(\frac{\partial F}{\partial s}(t,s)\frac{D\gamma'_{s}}{\D t}(t)\right)\D t\quad\mbox{by Fact~\eqref{eq:math151c:lecture12:fact1}}\\
&=\left(g_{\gamma_{s}(b)}\left(\frac{\partial F}{\partial s}(b,s),\gamma'_{s}(b)\right)-
g_{\gamma_{s}(a)}\left(\frac{\partial F}{\partial s}(a,s),\gamma'_{s}(a)\right)\right)\D t\nonumber\\
&\phantom{=\quad}-\int^{b}_{a}g_{\gamma_{s}(t)}\left(\frac{\partial F}{\partial s}(t,s)\frac{D\gamma'_{s}}{\D t}(t)\right)\D t\quad\mbox{by fundamental theorem of calculus},
  \end{align}
\end{subequations}
Then
\begin{equation}
\begin{split}
  \frac{\D}{\D s}E_{g}(\gamma_{s})&=g_{\gamma(b)}(V(b),\gamma'_{s}(b))
  -g_{\gamma(a)}(V(a),\gamma'_{s}(a))%
  \nonumber\\
&\phantom{=\quad}-\int^{b}_{a}g_{\gamma(t)}\left(V(t),\frac{D\gamma'}{\D t}(t)\right)\,\D t
\end{split}
\end{equation}
since $\partial_{s}F(t,s)|_{s=0}=V(t)$.
\end{proof}

\begin{remark}
Observe that metric-compatibility was used as a way to perform
``integration by parts''. This is how we use the condition in practice.
\end{remark}

\begin{corollary}
Let $(M,g)$ be a Riemannian manifold. We have $\gamma\colon[a,b]\to M$
is a geodesic if and only if it is a critical point of the energy,
i.e., for every valuation field $V\in\Vect{\gamma}$ we have $\D E_{g}(\gamma)[V]=0$.
\end{corollary}

\begin{theorem}
Let $(M,g)$ be a Riemannian manifold, let $\gamma\colon[a,b]\to M$ be
a smooth immersion. Then $\length_{g}$ is smooth [in the sense of
 Frechet derivatives of maps between Banach spaces] at $\gamma$ and
\begin{equation}
\D\length_{g}(\gamma)[V]
=\left.\frac{\D}{\D s}\length_{g}(\gamma_{s})\right|_{s=0}
=-\int^{b}_{a}g_{\gamma(t)}\left(V(t),\left(\frac{D}{\D t}\frac{\gamma'}{|\gamma'|_{g}}\right)(t)\right) \D t
\end{equation}
for every variation field $V$ along $\gamma$.
\end{theorem}

The proof is nearly identical to the previous theorem's proof.

\begin{definition}
Let $(M,g)$ be a Riemannian manifold, let $\gamma\colon[a,b]\to M$
be a smooth curve on $M$. We say $\gamma$ is a \define{Constant-Speed}
curve if $|\gamma'|_{g}$ is constant on $[a,b]$.
\end{definition}

\begin{corollary}
Let $(M,g)$ be a Riemannian manifold and let $\gamma\colon[a,b]\to M$
be a smooth immersion. Then $\gamma$ is a geodesic if and only if it
is a constant-speed critical point of the length.
\end{corollary}

\begin{proof}
\forwardproof\ Assume $\gamma$ is a geodesic. Then
\begin{subequations}
  \begin{align}
\frac{\D}{\D t}\left(|\gamma(t)|_{g}^{2}\right)
&=\frac{\D}{\D t}(g_{\gamma(t)}(\gamma'(t),\gamma'(t)))\\
&=2g_{\gamma(t)}\left(\frac{D\gamma'}{\D t}(t),\gamma'(t)\right)\\
&=2g_{\gamma(t)}\left(0,\gamma'(t)\right)=0.
  \end{align}
\end{subequations}
Therefore $|\gamma'|_{g}$ must be cosntant automatically. Then
\begin{equation}
\D\length_{g}(\gamma)[V]=\frac{1}{c}\D E_{g}(\gamma)[V]=0
\end{equation}
since $\gamma$ is a critical point of the energy.

\backwardproof\ Run the forward proof ``in reverse''.
\end{proof}

\begin{lemma}[Gauss]
Let $(M,g)$ be a Riemannian manifold, let $x\in M$, let $r\in(0,\operatorname{inj}_{x}(M))$.
For every tangent vector $v,w\in B_{r}(0)\subset T_{x}M$, we have
\begin{equation}
g_{\exp_{x}(v)}((\D\exp_{x})_{v}(v),(\D\exp_{x})_{v}(w))=g_{x}(v,w).
\end{equation}
\end{lemma}

\begin{proof}
Exercise.
\end{proof}

\begin{remark}
This is a sort of ``conformality'' result concerning the Gauss map.
\end{remark}